% ============================================================
%  SECTION 5: Appreciation and depreciation       (Member 2)
%  Lines marked [added at merge] are new; everything else is Member 2's text.
%  Fixed at merge: Ravi is Meera's uncle (not "whose uncle works in Dubai").
% ============================================================
\section{Appreciation and depreciation: who wins, who loses?}\label{sec:appdep}

A change in the exchange rate does not affect everyone in the same way. The result depends on whether a person or business needs to buy foreign currency to make payments or receives foreign currency as income.

\subsection{When the rupee depreciates}
Suppose the exchange rate rises from ₹82 to ₹86 per US dollar. The rupee has \keyterm{depreciated}{A fall in the value of a currency set by the market; here, more rupees are needed to buy one dollar.} because each US dollar now costs more rupees. This changes the cost of international transactions.

An Indian importer who pays a supplier in dollars must spend more rupees for the same dollar payment. The same applies to Meera, who is preparing to study in the United States: if her college fees are fixed in dollars, their rupee cost rises when the rupee depreciates.

For an Indian exporter, however, depreciation may be advantageous. Selvi runs a garment business in Tiruppur and receives payments from foreign buyers in dollars. When those dollars are converted into rupees, each dollar brings in more rupees. Indian goods may also become cheaper for foreign buyers when their prices are expressed in dollars, improving exporters' price competitiveness. The size of this advantage depends on how buyers respond to price changes.

Meera's uncle Ravi, who works in Dubai and sends money to India, illustrates another effect. If the remittance is sent in a foreign currency and its value against the rupee rises, the same amount of foreign currency converts into more rupees for the recipient.

\subsection{When the rupee appreciates}
If the exchange rate falls from ₹86 to ₹82 per US dollar, the rupee \keyterm{appreciates}{A rise in the value of a currency set by the market; here, fewer rupees are needed to buy one dollar.}. The effects are broadly reversed. Importers and students paying dollar-denominated expenses need fewer rupees to buy the same number of dollars. Meera's fees, for example, become less expensive in rupee terms.

Selvi may face the opposite effect: her dollar earnings convert into fewer rupees, and her garments may become more expensive for foreign buyers if their rupee prices remain unchanged. Ravi's family may likewise receive fewer rupees for the same foreign-currency remittance.

% [added at merge] devaluation and revaluation (used in Sections 6, 9 and 11)
Depreciation and appreciation describe a currency whose price is set by the market. When a government or central bank fixes the rate itself, a deliberate cut in the currency's official value is called a \keyterm{devaluation}{A deliberate cut in a currency's official value by the government or central bank under a fixed exchange rate.}, and a deliberate increase is a \keyterm{revaluation}{A deliberate increase in a currency's official value under a fixed exchange rate.}. India devalued the rupee in 1966 and again in 1991 \seesec{sec:systems}.

\subsection{Why the effects differ}
Table~\ref{tab:winners} summarises the main effects.

\begin{table}[!htb]
\centering\small\sffamily
\caption{Who gains and who loses when the rupee moves.}\label{tab:winners}
\begin{tabularx}{\linewidth}{@{}>{\raggedright\arraybackslash}p{33mm}>{\raggedright\arraybackslash}X>{\raggedright\arraybackslash}X@{}}
\toprule
\textbf{Group} & \textbf{Rupee depreciation} & \textbf{Rupee appreciation}\\
\midrule
Importers & Imports cost more in rupees & Imports cost less in rupees\\
Students paying in dollars & Expenses rise in rupee terms & Expenses fall in rupee terms\\
Indian exporters & May gain price competitiveness & May lose price competitiveness\\
Recipients of foreign-currency remittances & Receive more rupees for the same amount & Receive fewer rupees for the same amount\\
% [added at merge] three groups from the team plan
Foreign tourists in India & Holidays in India become cheaper & Holidays in India become dearer\\
Indian tourists abroad & Trips abroad cost more & Trips abroad cost less\\
Firms with dollar loans & Repayments cost more rupees & Repayments cost fewer rupees\\
\bottomrule
\end{tabularx}
\end{table}

Neither appreciation nor depreciation benefits everyone. A weaker rupee can help an exporter earning dollars while increasing the costs of an importer spending dollars. A stronger rupee can reduce the cost of foreign purchases while lowering the rupee value of foreign-currency earnings. The effect of an exchange-rate movement therefore depends on a person's or business's position in international transactions.
